\documentclass{article}
\usepackage{graphicx}
\usepackage{xcolor}
\usepackage{amsmath}
\usepackage{caption}
\usepackage{colortbl}

\definecolor{color1}{RGB}{200,20,200}
\colorlet{color2}{red!60!yellow}
\title{Assignment 4}
\author{Aryan Saraiya}
\begin{document}
	\maketitle
	
	\section{List}
	Following are some technologies that can be used to format content:
	\begin{itemize}
		\item \LaTeX
		\item HTML (+CSS)
		\item Markdown
	\end{itemize}
	
	\noindent You would most likely use \LaTeX \quad in the following scenarios:
	
	\begin{enumerate}
		\item To write research papers
		\item To embed math equations within other software (e.g. GDocs, HTML etc.)
		\item To create your CV, project reports \& presentations
	\end{enumerate}
	
	\noindent There are three main categories of electronic markup:
	
	\begin{description}
		\item[Presentational] Used by traditional word-processing systems: binary codes embedded
		within document text that produce the WYSIWYG (”what you see is what you
		get”) effect. e.g. Microsoft Word
		\item[Porcedural] Markup is embedded in text which provides instructions for programs to
		process the text. e.g. \TeX
		\item[Descriptive] Markup is specifically used to label parts of the document for what they
		are, rather than how they should be processed. e.g. \LaTeX, HTML5
	\end{description}
	
	\section{Tables \& Colors}
	
	Code for \textcolor{blue}{Table 1} is written in this subsection. However, it will be rendered on a
	special page. Change the optional arguments such that the table is drawn at the top
	or \fcolorbox{green}{white}{bottom} of a page or \fcolorbox{green}{white}{within} this section. The last row of this table shows the width
	of each column.
	
	\colorbox{orange}{Read more about how colors are represented as numbers.} Try to find out about the
	HTML notation (both hexadecimal and numeric). Define a custom color with 200 red,
	20 green and 200 blue. Define another custom color using colorlet with 60\% red, and
	40\% yellow. Now use these colors to write the following equation:
	
	\[
	\textcolor{color1}{T(n)} = \textcolor{color2}{aT(n/b) + f(n)}
	\]
	
	\section{Figures}
	\subsection{Pictures: rotation \& cropping}
	
	Include the MNIT logo and rotate it appropriately. Also include the MNIT building
	picture but crop the sky and make it is as wide as the text, as shown in Figure 2.
	
	\begin{figure}[h]
		\centering
		\includegraphics[scale=0.5]{Mnit_logo.png}
		\caption{MNIT logo}
	\end{figure}
	\begin{figure}[h]
		\centering
		\includegraphics[width=\textwidth, trim=0 0 0 150,clip]{Mnit_logo.png}
		\caption{This is a floating image and can appear anywhere.}
	\end{figure}
	
	
	\subsection{Multiple images}
	
	
	\begin{minipage}{0.45\textwidth}
		\includegraphics[width=\linewidth]{Mnit_logo.png}
		\captionof{figure}{This is the first minipage}
	\end{minipage}
	\quad
	\begin{minipage}{0.45\textwidth}
		\includegraphics[width=\linewidth]{Mnit_logo.png}
		\captionof{figure}{This is the second minipage}
	\end{minipage}
	
	
	\clearpage
	\begin{table}[h]
		\centering
		\begin{tabular}{|p{2.5cm}|p{2cm}|p{2cm}|p{3.5cm}|}
			\hline
			\rowcolor{green} Language & Requires compilation & Output & Use \\ \hline
			C/C++ & Yes & Executable & To make apps \\ \hline
			\LaTeX & Yes & PDF & To make documents \\ \hline
			HTML/CSS/JS & No & Webpage & To make websites \\ \hline
			\rowcolor{yellow} 2.5 cm & 2 cm & 2 cm & 3.5 cm \\ \hline
		\end{tabular}
		\caption{This table should be rendered on a special page.}
	\end{table}
	
	
\end{document}